\documentclass[12pt]{article}
\usepackage{graphicx} % Required for inserting images
\usepackage[portuguese]{babel}
\usepackage{url}
\usepackage[colorlinks=true,
    linkcolor=blue,
    citecolor=blue,
    urlcolor=blue
    ]{hyperref}
\usepackage{epigraph}
\usepackage{csquotes}
\usepackage[backend=biber,
    maxcitenames=3, 
    style=abnt,
    citestyle=abnt,
    date=year,
    language=brazilian,
    sorting=nyt,
    extrayear]{biblatex}
\usepackage{bookmark}
\usepackage[T1]{fontenc}
\usepackage[utf8]{inputenc}

\addbibresource{Khazzoom-Brookes.bib}


\title{Para uma crítica ontológica do postulado de Khazzoom-Brookes}
\author{Nicholas Funari Voltani}
\date{}

\begin{document}

\maketitle

\section{Contextualização do objeto da análise}
As crises do petróleo ao longo da década de 1970 instigaram discussões sobre medidas estatais para a estabilização de economias em seus âmbitos nacionais. Delas resultaram debates seminais na área da Economia da Energia: \textcite{Khazzoom1980,Khazzoom1987} e \textcite{Lovins1988} divergiram quanto à aderência de estimativas de poupança de energia de eletrodomésticos [\textit{appliances}] mais eficientes, e \textcite{Brookes1990,Brookes1992,Brookes1993} opôs-se a \textcite{Keepin1988,Keepin1989,Jones1989,Grubb1990,Grubb1992} quanto à aderência de aumentos de eficiência energética enquanto políticas para solucionar o aquecimento global (um problema que começava a ganhar tração política à época). 

O principal fruto destes debates foi a caracterização do efeito \textit{rebound}: o incentivo ao consumo \textit{aumentado} de energia após um aumento de eficiência energética, o que \textit{reduz} sua poupança esperada \parencite{Khazzoom1980,Brookes1990,Saunders2000,Saunders2009,Berkhout2000}. O postulado de Khazzoom-Brookes se trata de quando o efeito \textit{rebound} é intenso o suficiente para que a conservação energética \textit{esperada}, permitida pelo ganho de eficiência energética, seja totalmente neutralizada; dito de outra forma, se trata de quando a poupança \textit{efetiva} ou é nula ou mesmo \textit{negativa}, em que se consome \textit{mais} energia do que antes deste aumento de eficiência. Neste último caso, diz-se que houve um efeito \textit{backfire}. O postulado de Khazzoom-Brookes, por sua vez, assevera que, fixado o preço real da energia, o efeito \textit{rebound} ao nível da economia inteira [\textit{economy-wide}] tende a neutralizar totalmente a poupança esperada, podendo até mesmo levar a um efeito \textit{backfire} \parencite{Saunders1992,Saunders2000,Sorrell2009}.

\section{Problematização e conclusões da análise}
Os teóricos neoclássicos da Economia da Energia tendem a apreender o efeito \textit{rebound}/\textit{backfire} como se tratando de uma relação de causalidade do tipo ``se $X$, então $Y$''. Postulam, neste caso, que ganhos de eficiência energética são/podem ser \textit{causa} de aumentos de consumo energético, de tal forma que, caso tal relação seja verificada empiricamente, então ela estaria demonstrada como sendo verdadeira. Contudo, ao buscarem testar sua hipótese frente a situações concretas --- e.g. com dados econômicos da Suécia \parencite{Bergman1988}, da Escócia \parencite{Hanley2009} etc. ---, conseguem conclusões ambíguas, pois obtém-se tanto casos que confirmam tal relação, quanto casos que a negam \parencite{Greening2000,Sorrell2007}. 

O papel do presente trabalho é de subverter este raciocínio, buscando elucidar o vínculo necessário, no modo de produção capitalista, entre ganhos de eficiência energética e aumento de consumo energético. Valendo-se da teoria do valor-trabalho como exposta por Marx n'\textit{O capital} \parencite{Marx2017, Marx2014, Marx2017a}, assim como de trabalhos posteriores relacionados, como em \textcite{SaBarreto2018,SaBarreto2022}, apresenta-se a tese de que o postulado de Khazzoom-Brookes propõe uma relação de causalidade espúria, posto que tanto seu \textit{explanans} quanto seu \textit{explanandum} são consequências necessárias da própria lógica imanente do capital. 

Concentrando-se, a princípio, no nível de abstração da esfera da produção de capital --- dada sua \textit{prioridade ontológica} sobre a totalidade social \parencite{Lukacs2012,Lukacs2018} ---, é possível concluir que o imperativo do capital, que é a produção de mais-valor para fins de sua reprodução ampliada, engendra não só a escala crescente da esfera da produção, como também o revolucionamento contínuo das \textit{forças produtivas}; ou seja, engendra tanto o consumo crescente de meios de produção quanto o impulso ao seu uso mais ``produtivo''. Destaque-se, porém, que a categoria marxista de forças produtivas é mais ampla do que as categorias ``produtividade'' e ``eficiência'' da economia neoclásica. Na literatura marxista, \textit{produtividade} diz respeito à capacidade produtiva \textit{do trabalho}; \textit{eficiência}, dos \textit{meios de produção}. Porém, sob a perspectiva do capital --- e isso se reflete em sua representação neoclássica ---, ambos são indiferentes: importa-lhe produzir o mesmo produto com menos fatores de produção \parencite{SaBarreto2018,SaBarreto2022,Marx2017}. 


%Dessa forma, elucida-se que não é o aumento da eficiência \textit{per se} que induz o aumento do consumo energético, e sim que é \textit{a busca do capital pela valorização do valor} aquilo que engendra tanto o aumento do consumo de meios de produção --- dentre os quais fatores energéticos, como combustíveis ou eletricidade ---, como também a tendência ao aprimoramento das forças produtivas, o que inclui (mas não se esgota em) o aumento da eficiência dos processos produtivos no que diz respeito ao seu consumo de meios de produção. 

A principal conclusão obtida neste trabalho é que a tendência do consumo energético é \textit{crescente} em escala \textit{social}, conforme o sucesso do capital, considerado \textit{in totum}, em sua reprodução em escala ampliada. Ilustra-se, além disso, por que conclusões parciais sobre tais informações econômicas de países e/ou de setores industriais específicos podem aparecer fenomenicamente como o \textit{oposto} daquela tendência imanente do capital: porque, na realidade concreta, há inúmeras tendências concomitantemente em ação, reforçando-se e contrarrestando-se aqui e ali \parencite{Bhaskar1978,Bhaskar1998,Medeiros2013}.\footnote{\textcite[p. 15]{Sorrell2007} também reconhece isso, embora por motivos distintos, quando afirma que o escopo empírico do consumo energético deve ser o mais amplo possível.} O objetivo último deste trabalho é apresentar essa tendência fundamental, no cerne da sociedade capitalista, que inexoravelmente exerce influência, através da ``lei férrea'' da concorrência, sobre a esfera de ação do capitalista \textit{in actu} e, por consequência, sobre o planeta sob o jugo do capital. 

% A análise debruça-se especialmente sobre a esfera da produção --- assim como Marx o faz no livro I d'\textit{O capital} \parencite{Marx2017} ---, chegando à conclusão de que tanto o consumo energético quanto os níveis de eficiência energética são impulsionados, em última instância, pelo ímpeto intrínseco do capital de produção acelerada de mais-valor, que impele os capitais individuais, através da força da concorrência, ao melhoramento constante das forças produtivas sociais.

% A forma mais eminente pela qual o capital busca aumentar sua produção de mais-valor, além do aumento da intensidade do trabalho, é a substituição da força de trabalho por sistemas automatizados, em que o trabalhador assalariado passa agora a, como Marx o coloca, ``servir à máquina''. 

% Tal lógica é observável, por exemplo, no ímpeto do capital de aumentar não só a \textit{eficiência} no consumo produtivo de meios de produção --- requerer \textit{menos} destes para um nível de produção dado ---, mas também a \textit{produtividade} da produção --- em um dado período de tempo, produzir mais mercadorias, e, portanto, requerendo \textit{mais} meios de produção \parencite{SaBarreto2021a,SaBarreto2022}. Em suma: demonstra-se por que a ocorrência de efeitos \textit{backfire} do consumo energético, mediante aumentos da eficiência de seu uso, é uma \textit{lei de tendência} \parencite{Medeiros2013,Lawson1997,Bhaskar1978,Bhaskar1998} engendrada pela própria dinâmica da economia capitalista. 



\printbibliography


\end{document}
